\documentclass[aspectratio=169]{beamer}
\usetheme{Madrid}
\usepackage[spanish,es-nodecimaldot]{babel}
\usepackage[utf8]{inputenc}
\usepackage[T1]{fontenc}
\usepackage{amsmath}
\usepackage{amssymb}
\usepackage{graphicx}
\usepackage{booktabs}
\usepackage{xcolor}
\usepackage{verbatim}

% Color UNS
\definecolor{azulUNS}{RGB}{30,60,120}
\definecolor{grisclaro}{RGB}{245,245,245}
\usecolortheme[named=azulUNS]{structure}

\title[MAD1 -- U6]{Métodos de Análisis de Datos I}
\subtitle{Unidad 6 -- Bloque 5: Portfolio profesional}
\author{Departamento de Matemática -- UNS}
\date{2026}

\begin{document}

\begin{frame}
  \titlepage
\end{frame}

\begin{frame}{Contenidos}
  \tableofcontents
\end{frame}

%-----------------------------------------------
\section{¿Qué es un portfolio de data science?}
%-----------------------------------------------

\begin{frame}{¿Qué es un portfolio?}
  Un \textbf{portfolio} es una colección de proyectos que demuestra lo que sabés hacer con datos reales.

  \bigskip
  \begin{columns}
    \column{0.5\textwidth}
      \textbf{No es:}
      \begin{itemize}
        \item Una lista de cursos o certificados.
        \item Código copiado de tutoriales sin adaptación.
        \item Un proyecto con datos perfectos y resultados obvios.
        \item Una colección de notebooks sin explicar.
      \end{itemize}
    \column{0.5\textwidth}
      \textbf{Sí es:}
      \begin{itemize}
        \item Evidencia de que podés formular un problema, analizarlo y comunicar resultados.
        \item Trabajo propio, documentado, reproducible.
        \item Proyectos con datos reales o realistas, con decisiones justificadas.
        \item Algo que cualquiera pueda ejecutar y entender sin preguntarte.
      \end{itemize}
  \end{columns}

  \bigskip
  \begin{alertblock}{Principio}
    Un portfolio de 2 proyectos bien hechos vale más que 10 notebooks sin contexto.
  \end{alertblock}
\end{frame}

%-----------------------------------------------
\section{GitHub como vitrina}
%-----------------------------------------------

\begin{frame}{GitHub como vitrina profesional}
  GitHub es el estándar para mostrar trabajo técnico. Un perfil bien organizado comunica mucho antes de que alguien lea una línea de código.

  \bigskip
  \textbf{Buenas prácticas en GitHub:}
  \begin{itemize}
    \item Cada proyecto en un repositorio propio, con nombre descriptivo.
    \item \textbf{README.md} completo en cada repositorio (ver template más adelante).
    \item Commits con mensajes que explican qué se hizo, no solo ``update''.
    \item Código comentado: no en cada línea, pero sí en las decisiones no obvias.
    \item Sin archivos de datos sensibles ni credenciales subidas al repositorio.
    \item Notebooks con outputs limpios (sin errores ni celdas ejecutadas a medias).
  \end{itemize}

  \bigskip
  \begin{exampleblock}{Herramienta complementaria}
    Un blog personal o perfil en \textbf{Kaggle} / \textbf{Hugging Face} puede complementar GitHub para proyectos con visualizaciones interactivas o modelos deployados.
  \end{exampleblock}
\end{frame}

%-----------------------------------------------
\section{Estructura de un proyecto bien documentado}
%-----------------------------------------------

\begin{frame}{Estructura de carpetas recomendada}
  \begin{columns}
    \column{0.5\textwidth}
\texttt{mi\_proyecto/}\\
\texttt{├── README.md}\\
\texttt{├── data/}\\
\texttt{│\phantom{xx}├── raw/}\\
\texttt{│\phantom{xx}└── processed/}\\
\texttt{├── notebooks/}\\
\texttt{│\phantom{xx}├── 01\_exploracion.ipynb}\\
\texttt{│\phantom{xx}├── 02\_modelado.ipynb}\\
\texttt{│\phantom{xx}└── 03\_resultados.ipynb}\\
\texttt{├── src/}\\
\texttt{│\phantom{xx}└── utils.py}\\
\texttt{├── reports/}\\
\texttt{│\phantom{xx}└── reporte\_final.pdf}\\
\texttt{└── requirements.txt}
    \column{0.5\textwidth}
      \textbf{Notas:}
      \begin{itemize}
        \item \texttt{data/raw/}: datos originales, nunca modificados.
        \item \texttt{data/processed/}: datos limpios listos para modelar.
        \item Los notebooks se numeran para indicar el orden de ejecución.
        \item \texttt{src/} para funciones reutilizables (no repetir código entre notebooks).
        \item \texttt{requirements.txt} permite reproducir el entorno.
      \end{itemize}
      \bigskip
      \begin{alertblock}{Regla}
        Si los datos son grandes o sensibles, subir solo una muestra o instrucciones para obtenerlos.
      \end{alertblock}
  \end{columns}
\end{frame}

\begin{frame}{Template: README.md de un proyecto}
  \textbf{Secciones mínimas de un buen README:}

  \bigskip
  \begin{enumerate}
    \item \textbf{Título y descripción breve} (2--3 líneas): qué problema resuelve y con qué datos.
    \item \textbf{Motivación}: por qué el problema es interesante o relevante.
    \item \textbf{Datos}: origen, tamaño, variables principales.
    \item \textbf{Metodología}: qué técnicas se usaron y por qué.
    \item \textbf{Resultados principales}: métrica clave y hallazgo más relevante en 1--2 oraciones.
    \item \textbf{Cómo ejecutar}: instrucciones claras para reproducir el análisis.
    \item \textbf{Limitaciones y próximos pasos}.
  \end{enumerate}

  \bigskip
  \begin{exampleblock}{Criterio simple}
    Si alguien ajeno al proyecto puede entender qué hiciste, por qué y qué encontraste leyendo solo el README, está bien escrito.
  \end{exampleblock}
\end{frame}

%-----------------------------------------------
\section{Qué mostrar de cada unidad}
%-----------------------------------------------

\begin{frame}{Qué podés mostrar: Unidad 1 -- Introducción y EDA}
  \textbf{Proyecto sugerido:} análisis exploratorio completo de un dataset de interés personal.

  \bigskip
  \textbf{Lo que demuestra:}
  \begin{itemize}
    \item Capacidad de limpiar datos reales (valores faltantes, outliers, tipos).
    \item Exploración visual con criterio: no graficar todo, sino lo que responde preguntas.
    \item Formulación de hipótesis a partir de los datos.
  \end{itemize}

  \bigskip
  \textbf{Mini template de notebook:}
  \begin{enumerate}
    \item \textbf{Celda 1}: contexto del dataset (origen, tamaño, pregunta que guía el EDA).
    \item \textbf{Celda 2}: carga y revisión inicial (\texttt{head}, \texttt{info}, \texttt{describe}).
    \item \textbf{Celda 3}: tratamiento de valores faltantes con justificación.
    \item \textbf{Celda 4}: distribuciones de variables clave (histogramas, boxplots).
    \item \textbf{Celda 5}: relaciones entre variables (correlación, dispersión).
    \item \textbf{Celda 6}: síntesis en texto: qué encontré y qué pregunta surge.
  \end{enumerate}
\end{frame}

\begin{frame}{Qué podés mostrar: Unidad 2 -- Regresión}
  \textbf{Proyecto sugerido:} predicción de una variable continua con interpretación de coeficientes.

  \bigskip
  \textbf{Lo que demuestra:}
  \begin{itemize}
    \item Construcción y evaluación de un modelo de regresión lineal o logística.
    \item Interpretación del significado de cada coeficiente en términos del dominio.
    \item Verificación de supuestos (linealidad, homocedasticidad, normalidad de residuos).
  \end{itemize}

  \bigskip
  \textbf{Mini template de notebook:}
  \begin{enumerate}
    \item \textbf{Celda 1}: formulación del problema y justificación del modelo.
    \item \textbf{Celda 2}: preprocesamiento (encoding, escalado si corresponde).
    \item \textbf{Celda 3}: ajuste del modelo e interpretación de coeficientes.
    \item \textbf{Celda 4}: gráficos de residuos y verificación de supuestos.
    \item \textbf{Celda 5}: métricas ($R^2$, RMSE, MAE) en train y test.
    \item \textbf{Celda 6}: conclusión: ¿qué variables importan y cuánto?
  \end{enumerate}
\end{frame}

\begin{frame}{Qué podés mostrar: Unidad 3 -- Clasificación}
  \textbf{Proyecto sugerido:} clasificador binario o multiclase con evaluación completa.

  \bigskip
  \textbf{Lo que demuestra:}
  \begin{itemize}
    \item Manejo del desbalance de clases y elección del umbral de decisión.
    \item Uso adecuado de métricas (precisión, recall, F1, AUC-ROC) según el contexto.
    \item Comparación entre al menos dos modelos con justificación de la elección final.
  \end{itemize}

  \bigskip
  \textbf{Mini template de notebook:}
  \begin{enumerate}
    \item \textbf{Celda 1}: descripción del problema y análisis del desbalance de clases.
    \item \textbf{Celda 2}: pipeline de preprocesamiento con \texttt{Pipeline} de sklearn.
    \item \textbf{Celda 3}: entrenamiento de dos o más modelos con validación cruzada.
    \item \textbf{Celda 4}: curva ROC, matriz de confusión, reporte de clasificación.
    \item \textbf{Celda 5}: análisis del umbral: trade-off precisión/recall.
    \item \textbf{Celda 6}: justificación del modelo final y sus limitaciones.
  \end{enumerate}
\end{frame}

\begin{frame}{Qué podés mostrar: Unidad 4 -- Modelos avanzados}
  \textbf{Proyecto sugerido:} comparación de modelos de ensamble o SVM con optimización de hiperparámetros.

  \bigskip
  \textbf{Lo que demuestra:}
  \begin{itemize}
    \item Uso de Random Forest, Gradient Boosting o SVM en un problema real.
    \item Búsqueda sistemática de hiperparámetros (GridSearch, RandomSearch, Optuna).
    \item Análisis de importancia de variables e interpretación del modelo.
  \end{itemize}

  \bigskip
  \textbf{Mini template de notebook:}
  \begin{enumerate}
    \item \textbf{Celda 1}: justificación de por qué un modelo complejo sobre uno simple.
    \item \textbf{Celda 2}: línea de base con modelo simple (regresión logística, árbol).
    \item \textbf{Celda 3}: entrenamiento del modelo avanzado con hiperparámetros por defecto.
    \item \textbf{Celda 4}: optimización de hiperparámetros y curva de validación.
    \item \textbf{Celda 5}: importancia de variables y análisis de errores.
    \item \textbf{Celda 6}: comparación con la línea de base: ¿valió la complejidad?
  \end{enumerate}
\end{frame}

\begin{frame}{Qué podés mostrar: Unidad 5 -- Aprendizaje no supervisado}
  \textbf{Proyecto sugerido:} segmentación de entidades con clustering e interpretación de grupos.

  \bigskip
  \textbf{Lo que demuestra:}
  \begin{itemize}
    \item Aplicación de clustering (K-Means, DBSCAN, jerárquico) con selección justificada del número de grupos.
    \item Reducción de dimensionalidad (PCA, UMAP) para visualización.
    \item Interpretación de los clusters en términos del dominio.
  \end{itemize}

  \bigskip
  \textbf{Mini template de notebook:}
  \begin{enumerate}
    \item \textbf{Celda 1}: motivación para segmentar (¿qué decisión se toma con los grupos?).
    \item \textbf{Celda 2}: preprocesamiento y escalado (crítico en clustering).
    \item \textbf{Celda 3}: selección del número de clusters (codo, silueta, dendrograma).
    \item \textbf{Celda 4}: visualización en 2D con PCA o UMAP.
    \item \textbf{Celda 5}: perfil de cada cluster (estadísticas descriptivas por grupo).
    \item \textbf{Celda 6}: interpretación: ¿qué nombre le pondrías a cada grupo y por qué?
  \end{enumerate}
\end{frame}

%-----------------------------------------------
\section{Errores comunes en portfolios}
%-----------------------------------------------

\begin{frame}{Errores comunes en portfolios}
  \begin{columns}
    \column{0.5\textwidth}
      \textbf{En el código:}
      \begin{itemize}
        \item Notebooks con celdas desordenadas o sin ejecutar de arriba a abajo.
        \item Variables con nombres \texttt{df2}, \texttt{X\_final2}, \texttt{model\_v3}.
        \item Sin separación entre exploración, modelado y evaluación.
        \item Código repetido que debería estar en funciones.
      \end{itemize}
    \column{0.5\textwidth}
      \textbf{En el contenido:}
      \begin{itemize}
        \item Métricas en entrenamiento sin reportar test.
        \item Ninguna discusión de limitaciones.
        \item Conclusiones que no se desprenden del análisis.
        \item Dataset de Titanic, Iris o tips sin ninguna vuelta de tuerca propia.
        \item Copiar código de Kaggle sin entenderlo ni adaptarlo.
      \end{itemize}
  \end{columns}

  \bigskip
  \begin{alertblock}{La señal más clara de un buen proyecto}
    El analista tomó decisiones y las justificó. No siguió un tutorial: resolvió un problema.
  \end{alertblock}
\end{frame}

%-----------------------------------------------
\section{Cierre}
%-----------------------------------------------

\begin{frame}{Ideas para llevarse}
  \begin{enumerate}
    \item Un portfolio demuestra proceso, no solo resultados.
    \item GitHub es la vitrina estándar: mantenerlo ordenado y documentado.
    \item Cada proyecto necesita un README que permita entenderlo sin preguntar.
    \item Cada unidad de la materia es base para un proyecto de portfolio real.
    \item La estructura del notebook importa: contexto $\rightarrow$ datos $\rightarrow$ análisis $\rightarrow$ conclusión.
    \item Evitar los datasets cliché sin adaptación propia.
    \item La señal de un buen analista: tomó decisiones y las justificó.
  \end{enumerate}

  \bigskip
  \begin{exampleblock}{Para empezar hoy}
    Elegir un dataset que te interese de verdad. Formular una pregunta concreta. Seguir el template de la unidad más cercana al problema. Documentar cada decisión. Subir a GitHub.
  \end{exampleblock}
\end{frame}

\end{document}
